\section{数据表示、内存、指针和对象布局}

\subsection{本章要解决的问题：\code{a[i]} 到底访问了什么}

刚学 C/C++ 时，数组访问通常被解释成“取第 \code{i} 个元素”。这个说法对写程序有用，但对理解机器执行和性能远远不够。高性能计算和 AI CPU 算子优化中，大量核心问题都可以追到一个更底层的问题：

\begin{lstlisting}
某条指令到底访问了哪个地址上的哪些字节？
\end{lstlisting}

先看一个简单函数：

\begin{lstlisting}[language=C++]
int get(int const* a, int i)
{
    return a[i];
}
\end{lstlisting}

在 x86-64 System V ABI 下，编译器可能生成类似汇编：

\begin{lstlisting}
movsxd  rsi, esi
mov     eax, dword ptr [rdi + rsi*4]
ret
\end{lstlisting}

这三行非常值得细看。

\begin{itemize}
  \item \register{rdi} 保存第一个参数 \code{a}，也就是数组首元素地址。
  \item \register{esi} 保存第二个参数 \code{i} 的低 32 位。
  \item \instruction{movsxd rsi, esi} 把 32 位有符号整数扩展成 64 位索引。
  \item \code{[rdi + rsi*4]} 是地址表达式，意思是 \code{a + i * 4}。
  \item \code{dword ptr} 说明这次从内存读取 4 字节。
  \item \register{eax} 是返回值寄存器的低 32 位。
\end{itemize}

于是 \code{a[i]} 在机器层面不是一个神秘的“数组操作”，而是一次地址计算加一次内存读取：

\begin{lstlisting}
address = base_address(a) + i * sizeof(int)
load 4 bytes from address
interpret those 4 bytes as int
\end{lstlisting}

\begin{keyidea}
机器没有“数组第几个元素”的概念。机器看到的是地址、字节数、寄存器和指令。C/C++ 类型系统把“第几个元素”翻译成“从哪个基地址开始，按多大步长计算地址，读写多少字节，并按什么规则解释这些字节”。
\end{keyidea}

本章要把这个过程讲清楚。后面所有性能优化几乎都会反复依赖本章内容：cache line 为什么一次拉 64 字节，SIMD load 为什么要求连续布局，结构体字段顺序为什么影响性能，未定义行为为什么会让编译器大胆优化，AoS 和 SoA 为什么差异巨大，GEMM packing 为什么本质上是重排内存访问。

\subsection{本章学习路线：从一个字节走到一个 kernel}

本章容易学散，因为它同时出现位、整数、指针、数组、结构体、padding、cache line 和汇编寻址。不要把它们当成孤立知识点背。它们之间有一条非常明确的依赖链：

\begin{lstlisting}
bit -> byte -> object representation -> type rule
    -> address -> pointer expression -> layout
    -> x86-64 addressing mode -> cache line behavior
    -> SIMD / HPC / AI kernel data movement
\end{lstlisting}

这条链路的意思是：

\begin{enumerate}
  \item \textbf{bit 和 byte} 让你知道机器保存信息的最小粒度。
  \item \textbf{对象表示} 让你知道一个 C++ 对象最终落成哪些字节。
  \item \textbf{类型规则} 让你知道这些字节如何被解释，编译器能假设什么。
  \item \textbf{地址和指针} 让你把源码表达式变成具体内存位置。
  \item \textbf{布局} 让你知道数组元素、结构体字段、padding 在地址空间里的相对位置。
  \item \textbf{x86-64 寻址模式} 让你把 C++ 访问表达式对应到 \code{[base + index * scale + displacement]}。
  \item \textbf{cache line} 让你解释为什么同样的运算量会因为访问模式不同而差很多。
  \item \textbf{SIMD 和 kernel 优化} 则要求你主动设计数据布局，让硬件以更少指令、更少 cache miss、更高带宽效率完成同样计算。
\end{enumerate}

如果你只是记住“int 通常 4 字节”“double 通常 8 字节”，本章价值很小。真正要训练的是一种固定能力：看到任意 C++ 表达式，能回答它需要访问哪些字节；看到任意内存操作汇编，能反推出它可能来自什么数据布局；看到一个慢循环，能判断它是不是被布局、步长、对齐或别名拖慢。

\subsubsection{本章的四个观察层次}

同一个程序可以从四个层次观察。初学者常犯的问题，是只停在第一层或把四层混在一起。

\begin{center}
\begin{tabular}{p{0.16\linewidth}p{0.36\linewidth}p{0.36\linewidth}}
\toprule
层次 & 你看到什么 & 本章要你学会的问题 \\
\midrule
C++ 源码层 & \code{a[i]}、\code{p->x}、\code{sizeof(S)} & 这个表达式的类型是什么？对象生命周期存在吗？访问是否越界？ \\
对象布局层 & 字段 offset、padding、对齐、数组步长 & 目标字节从基地址偏移多少？一次访问读写几字节？ \\
汇编层 & \code{mov eax, [rdi+rsi*4]} & 基址、索引、缩放和常量偏移分别来自源码里的什么？ \\
微架构预告层 & cache line、步长、连续性、带宽 & 这些访问会顺序命中 cache，还是浪费 cache line、破坏预取和 SIMD？ \\
\bottomrule
\end{tabular}
\end{center}

学习时必须强迫自己每次都跨层解释。例如：

\begin{lstlisting}[language=C++]
float z = particles[i].z;
\end{lstlisting}

不要只说“取第 i 个粒子的 z”。更完整的解释应该是：

\begin{enumerate}
  \item \code{particles} 是 \code{Particle*}，所以 \code{particles + i} 以 \code{sizeof(Particle)} 为步长。
  \item 如果 \code{Particle} 大小为 16，\code{z} 字段 offset 为 8，那么目标地址是“首元素地址 + \code{i * 16 + 8}”。
  \item 目标字段是 \code{float}，因此 load 宽度是 4 字节。
  \item 在 x86-64 汇编里，因为 scale 不能直接取 16，编译器可能先计算 \code{i * 16}，再访问 \code{dword ptr [base + offset + 8]}。
  \item 如果循环只读 \code{z}，AoS 布局中每 16 字节只用 4 字节，cache line 利用率可能只有四分之一。
\end{enumerate}

对应的汇编形态可以放进单独代码块观察：

\begin{lstlisting}
shl   rax, 4
movss xmm0, dword ptr [rdi + rax + 8]
\end{lstlisting}

这种解释方式看起来繁琐，但它是性能优化的基本功。后面做矩阵乘、卷积、attention、quantization、packing、prefetch、SIMD load/store 时，所有复杂技术都只是这条链路的规模化版本。

\subsubsection{符号约定和推理习惯}

为了让后面的推导不混乱，本章使用下面几个约定：

\begin{itemize}
  \item \code{addr(x)} 表示对象或子对象 \code{x} 的起始地址。
  \item \code{sizeof(T)} 表示类型 \code{T} 的对象大小，单位是字节。
  \item \code{offset(S, field)} 表示字段 \code{field} 相对结构体 \code{S} 起始地址的字节偏移。
  \item \code{stride} 表示相邻两个逻辑元素之间的地址距离，单位通常是字节。
  \item \code{load width} 表示一次 load 指令读取多少字节。
\end{itemize}

你每次推导内存访问时，都应该把“元素下标”和“字节偏移”分开写。下标是源码层概念，偏移是机器地址层概念。例如：

\begin{lstlisting}
a[3] index        = 3
a[3] byte_offset  = 3 * sizeof(int)
                  = 12
\end{lstlisting}

把这两个数混在一起，是阅读汇编时最常见的早期错误。

\begin{exercisebox}
\textbf{随堂检查 3.0：把一句话说完整}

请不要查资料，直接回答下面三个问题。每个答案至少写出“类型、基地址、字节偏移、访问宽度”四个词。

\begin{enumerate}
  \item \code{double* p; p[i]} 访问的字节地址如何计算？
  \item \code{struct V \{ float x, y, z, w; \}; V* v; v[i].z} 访问的字节地址如何计算？
  \item 下面两条汇编合起来可能对应什么 C++ 访问？
\end{enumerate}

\begin{lstlisting}
shl   rax, 4
movss xmm0, dword ptr [rdi + rax + 8]
\end{lstlisting}

如果你的答案里只有“取第几个元素”，说明你还停留在源码层；本章后面要反复训练到机器层。
\end{exercisebox}

\subsection{位、字节和机器字：数据表示的最小地基}

\subsubsection{位是信息，字节是寻址单位}

\term{bit}{binary digit} 是二进制位，只能表示 0 或 1。一个 bit 可以表达一个真假状态，但真实程序中的整数、地址、浮点数和指令都需要多个 bit 组合。

\term{byte}{字节} 在现代通用机器上通常是 8 bit。C++ 标准中 \code{char} 是最小可寻址对象单位；在我们学习的 x86-64/Linux 环境中，一个 byte 就是 8 bit。

\begin{center}
\begin{tabular}{lll}
\toprule
单位 & 大小 & 常见用途 \\
\midrule
bit & 1 个二进制位 & 标志位、掩码的一位 \\
byte & 8 bit & 最小寻址单位、\code{char}、对象表示 \\
word & 依上下文变化 & CPU 自然处理宽度或 ISA 术语 \\
dword & 32 bit & x86 汇编中的双字 \\
qword & 64 bit & x86 汇编中的四字 \\
\bottomrule
\end{tabular}
\end{center}

注意 \term{word}{字} 是一个容易混淆的词。在 x86 历史语境里，word 常指 16 bit，dword 是 32 bit，qword 是 64 bit；在体系结构教材中，word 有时指机器自然处理宽度；在内存系统里，cache line 又是另一种粒度。读文档时必须看上下文。

\subsubsection{二进制、十六进制和为什么程序员爱用十六进制}

二进制直接对应 bit，但写起来太长。十六进制每一位刚好表示 4 bit，因此两个十六进制数字刚好表示 1 byte：

\begin{center}
\begin{tabular}{lll}
\toprule
十六进制 & 二进制 & 十进制 \\
\midrule
\code{0x0} & \code{0000} & 0 \\
\code{0x1} & \code{0001} & 1 \\
\code{0x7} & \code{0111} & 7 \\
\code{0x8} & \code{1000} & 8 \\
\code{0xF} & \code{1111} & 15 \\
\bottomrule
\end{tabular}
\end{center}

例如：

\begin{lstlisting}
0x12 = 0001 0010
0x34 = 0011 0100
0x1234 = 0001 0010 0011 0100
\end{lstlisting}

十六进制非常适合观察内存、机器码和地址。后面你会经常看到：

\begin{itemize}
  \item 对象字节：\code{78 56 34 12}。
  \item 机器码字节：\code{8b 04 b7}。
  \item 地址：\code{0x7ffe...}。
  \item 位掩码：\code{0xff}、\code{0xffff0000}。
\end{itemize}

\subsection{内存的第一模型：按字节寻址的巨大数组}

从最简单的机器模型看，内存可以被想象成一个按地址索引的字节数组：

\begin{lstlisting}
address    byte
0x1000  -> 0x78
0x1001  -> 0x56
0x1002  -> 0x34
0x1003  -> 0x12
0x1004  -> ...
\end{lstlisting}

地址是“某个字节的位置编号”。如果一个 \code{int} 占 4 字节，地址 \code{0x1000} 处的 \code{int} 会覆盖：

\begin{lstlisting}
0x1000
0x1001
0x1002
0x1003
\end{lstlisting}

但是这个模型还不完整，因为普通 Linux 进程看到的是虚拟地址，不是物理内存条上的真实位置。更准确的说法是：对用户态程序而言，指针值通常是虚拟地址；CPU 执行 load/store 时会通过地址翻译机制把虚拟地址转换为物理地址，期间会用到 TLB 和页表。

本章先使用“内存是字节数组”的模型来理解对象表示、指针和布局。第 7 章、第 33 章会把虚拟内存、页表和 TLB 展开。

\begin{mentalmodel}
先把内存想成字节数组，但永远记住它只是第一模型。性能优化时，你还要逐步叠加 cache line、cache set、TLB、page、NUMA、预取器和内存带宽。
\end{mentalmodel}

\subsection{类型不是字节本身，而是解释字节的规则}

内存中的 byte 本身没有“这是 int”或“这是 float”的标签。类型存在于程序语言、编译器和指令语义中。

例如同样 4 个字节：

\begin{center}
\begin{minipage}{0.72\linewidth}
\begin{verbatim}
00 00 80 3f
\end{verbatim}
\end{minipage}
\end{center}

如果按 32 位整数解释，它是某个整数位模式；如果按 IEEE 754 \code{float} 解释，它是 \code{1.0f}；如果按 4 个 \code{unsigned char} 解释，它只是 4 个字节。

C++ 的类型系统至少做了四件事：

\begin{enumerate}
  \item 决定对象需要多少字节，例如 \code{sizeof(int)}。
  \item 决定对象要求什么对齐，例如 \code{alignof(double)}。
  \item 决定表达式如何生成地址，例如 \code{p + 1} 前进多少字节。
  \item 决定编译器可以做哪些优化，例如别名规则和对象生命周期规则。
\end{enumerate}

\begin{warningbox}
“内存只是字节，所以我想怎么解释就怎么解释”是非常危险的半真半假。硬件层面确实是字节；C++ 语言层面有对象、类型、生命周期、对齐和别名规则。高性能代码经常需要看字节，但必须用语言允许的方式看。
\end{warningbox}

\subsection{对象表示：C++ 对象在内存中的字节}

C++ 中的对象有类型、存储期、生命周期和值。对象在内存中的底层字节叫作 \term{object representation}{对象表示}。对于一个可平凡复制的对象，你可以通过 \code{unsigned char}、\code{std::byte} 或 \code{char} 观察它的对象表示。

下面的程序打印一个 32 位整数的字节：

\begin{lstlisting}[language=C++]
#include <cstdint>
#include <iomanip>
#include <iostream>

int main()
{
    std::uint32_t x = 0x12345678U;
    auto const* p = reinterpret_cast<unsigned char const*>(&x);

    for (std::size_t i = 0; i < sizeof(x); ++i) {
        std::cout << std::hex << std::setw(2) << std::setfill('0')
                  << static_cast<unsigned>(p[i]) << '\n';
    }
}
\end{lstlisting}

在 x86-64 上，你通常会看到：

\begin{center}
\begin{minipage}{0.5\linewidth}
\begin{verbatim}
78
56
34
12
\end{verbatim}
\end{minipage}
\end{center}

这就是小端字节序。

\subsection{字节序：为什么 \code{0x12345678} 在内存里像反过来}

\term{endianness}{字节序} 描述多字节对象的字节如何放进连续地址。

\subsubsection{小端和大端}

如果最低有效字节放在最低地址，叫 \term{little-endian}{小端}。x86-64 使用小端。对 \code{0x12345678}：

\begin{center}
\begin{tabular}{cc}
\toprule
地址偏移 & 字节 \\
\midrule
\code{+0} & \code{0x78} \\
\code{+1} & \code{0x56} \\
\code{+2} & \code{0x34} \\
\code{+3} & \code{0x12} \\
\bottomrule
\end{tabular}
\end{center}

如果最高有效字节放在最低地址，叫 \term{big-endian}{大端}。同一个值会存成：

\begin{center}
\begin{tabular}{cc}
\toprule
地址偏移 & 字节 \\
\midrule
\code{+0} & \code{0x12} \\
\code{+1} & \code{0x34} \\
\code{+2} & \code{0x56} \\
\code{+3} & \code{0x78} \\
\bottomrule
\end{tabular}
\end{center}

\subsubsection{字节序影响什么，不影响什么}

字节序影响的是多字节对象在内存中的排列。它不影响整数的数学值，也不影响寄存器里 \code{0x12345678} 这个值本身。

字节序经常在这些场景中变重要：

\begin{itemize}
  \item 打印对象字节。
  \item 解析二进制文件格式。
  \item 网络协议和磁盘格式。
  \item 手写 SIMD shuffle 或处理 packed data。
  \item 调试机器码和内存 dump。
\end{itemize}

多数纯 C++ 算法不需要直接关心字节序，但系统性能工程师必须能看懂字节序，否则看到内存 dump 时会误读数据。

\subsection{整数表示：无符号模运算和有符号补码}

\subsubsection{无符号整数是按位宽取模的环}

一个 N 位无符号整数可以表示 $0$ 到 $2^N - 1$。运算按模 $2^N$ 回绕。例如 8 位无符号整数：

\begin{center}
\begin{minipage}{0.6\linewidth}
\begin{verbatim}
255 + 1 == 0
0 - 1 == 255
\end{verbatim}
\end{minipage}
\end{center}

C++ 对无符号整数溢出有明确规定：按模回绕。这让无符号整数适合写位运算、哈希、掩码、循环计数的一些底层逻辑。

\subsubsection{有符号整数通常用补码表示}

现代 x86-64 使用 \term{two's complement}{二进制补码} 表示有符号整数。以 8 位为例：

\begin{center}
\begin{tabular}{lll}
\toprule
位模式 & 无符号解释 & 有符号补码解释 \\
\midrule
\code{0000 0000} & 0 & 0 \\
\code{0000 0001} & 1 & 1 \\
\code{0111 1111} & 127 & 127 \\
\code{1000 0000} & 128 & -128 \\
\code{1111 1111} & 255 & -1 \\
\bottomrule
\end{tabular}
\end{center}

补码的好处是加法电路可以同时服务有符号和无符号加法。比如 \code{-1} 的 8 位补码是 \code{1111 1111}，加 1 得到 \code{0000 0000}，自然回到 0。

\subsubsection{语言规则和机器行为不能混为一谈}

机器补码加法会回绕，但 C++ 有符号整数溢出是未定义行为。也就是说，下面代码在 C++ 语言层面不能被解释成“最大值加一变最小值”：

\begin{lstlisting}[language=C++]
int f(int x)
{
    return x + 1;
}
\end{lstlisting}

如果 \code{x == INT_MAX}，表达式 \code{x + 1} 发生有符号溢出，程序有未定义行为。编译器可以假设这种情况不会发生，并据此优化。

这不是抠字眼。性能优化中，编译器对循环边界、条件判断、向量化和死代码删除的很多决定，都依赖“程序没有未定义行为”的假设。

\begin{deepdive}
后续学习编译器优化时，你会看到 signed overflow、strict aliasing、out-of-bounds、未初始化读取、对象生命周期等规则如何影响优化。底层工程不是绕开语言规则，而是理解语言规则如何给编译器提供证明空间。
\end{deepdive}

\subsection{浮点数只先建立直觉}

本章重点是整数、地址和对象布局。浮点数会在 AI 算子、数值稳定性、SIMD 和 FMA 章节中深入。本章先建立三个直觉：

\begin{enumerate}
  \item \code{float} 和 \code{double} 也是固定字节数的对象表示。
  \item 它们通常使用 IEEE 754 格式，包含符号、指数和尾数。
  \item 浮点运算不是实数运算，存在舍入、特殊值、非结合性和 fast-math 语义问题。
\end{enumerate}

例如 \code{float x = 1.0f;} 的对象字节在小端 x86-64 上通常是：

\begin{center}
\begin{minipage}{0.5\linewidth}
\begin{verbatim}
00 00 80 3f
\end{verbatim}
\end{minipage}
\end{center}

按整数看这只是一个位模式；按 IEEE 754 \code{float} 看它表示 1.0。这个例子再次提醒你：字节没有类型，解释规则来自类型和上下文。

\subsection{指针：地址值加类型语义}

\subsubsection{指针值通常是地址}

在普通 x86-64/Linux 用户态程序中，指针值通常可以理解为虚拟地址。比如：

\begin{lstlisting}[language=C++]
int x = 42;
int* p = &x;
\end{lstlisting}

\code{p} 保存的是对象 \code{x} 的地址。打印 \code{p} 可能得到类似：

\begin{center}
\begin{minipage}{0.62\linewidth}
\begin{verbatim}
0x7ffdf2c3a6bc
\end{verbatim}
\end{minipage}
\end{center}

但指针不仅是一个整数地址。C++ 指针还携带类型语义：\code{int*} 指向 \code{int} 对象，\code{double*} 指向 \code{double} 对象，\code{Particle*} 指向 \code{Particle} 对象。类型决定了解引用读取多少字节、指针加法前进多少字节、编译器如何判断别名。

\subsubsection{解引用是一次按类型解释的内存访问}

表达式 \code{*p} 不是“取地址里的东西”这么笼统。更准确地说：

\begin{enumerate}
  \item 读取指针表达式 \code{p} 的值，得到一个地址。
  \item 根据 \code{p} 的静态类型知道目标对象类型是 \code{int}。
  \item 从该地址开始读取 \code{sizeof(int)} 个字节。
  \item 按 \code{int} 的对象表示解释这些字节。
\end{enumerate}

如果地址没有对齐、对象生命周期不存在、类型不匹配或越界，语言层面可能已经出错。硬件是否“碰巧能读”不是 C++ 程序正确性的充分条件。

\subsubsection{空指针、野指针和悬垂指针}

指针错误是底层学习必须严肃对待的内容。

\begin{description}
  \item[空指针] 不指向任何对象。解引用空指针是错误。
  \item[野指针] 未初始化或保存无效地址的指针。它可能指向任意位置。
  \item[悬垂指针] 原来指向有效对象，但对象生命周期已经结束。
\end{description}

例如：

\begin{lstlisting}[language=C++]
int* bad()
{
    int x = 1;
    return &x;
}
\end{lstlisting}

\code{x} 是局部自动对象，函数返回后生命周期结束。返回的指针值可能仍然看起来像一个地址，但它不再指向活着的 \code{int} 对象。用它做性能实验，结果没有意义。

\subsection{指针加法：类型决定缩放}

指针加法不是简单整数加法。若 \code{p} 类型是 \code{T*}，则 \code{p + k} 的地址通常是：

\begin{center}
\begin{minipage}{0.76\linewidth}
\begin{verbatim}
address(p + k) = address(p) + k * sizeof(T)
\end{verbatim}
\end{minipage}
\end{center}

例如：

\begin{lstlisting}[language=C++]
int* p = reinterpret_cast<int*>(0x1000);
int* q = p + 3;
\end{lstlisting}

如果 \code{sizeof(int) == 4}，则 \code{q} 的数值地址是 \code{0x100c}，不是 \code{0x1003}。

这解释了 \code{a[i]}：

\begin{center}
\begin{minipage}{0.7\linewidth}
\begin{verbatim}
a[i] == *(a + i)
\end{verbatim}
\end{minipage}
\end{center}

如果 \code{a} 是 \code{int*}，\code{a + i} 自动按 4 缩放；如果 \code{a} 是 \code{double*}，按 8 缩放；如果 \code{a} 是 \code{Particle*}，按 \code{sizeof(Particle)} 缩放。

\begin{warningbox}
\code{reinterpret_cast<std::uintptr_t>(p) + 1} 和 \code{p + 1} 不是同一个概念。前者是整数地址加 1 字节；后者是指针按元素大小前进 1 个对象。写底层代码时必须分清“字节偏移”和“元素偏移”。
\end{warningbox}

\subsection{数组：连续对象序列}

C/C++ 数组是连续存放的同类型对象序列：

\begin{lstlisting}[language=C++]
int a[4] = {10, 20, 30, 40};
\end{lstlisting}

如果 \code{a[0]} 地址为 \code{0x1000}，且 \code{sizeof(int) == 4}，则布局通常是：

\begin{center}
\begin{tabular}{lll}
\toprule
元素 & 起始地址 & 覆盖字节范围 \\
\midrule
\code{a[0]} & \code{0x1000} & \code{0x1000..0x1003} \\
\code{a[1]} & \code{0x1004} & \code{0x1004..0x1007} \\
\code{a[2]} & \code{0x1008} & \code{0x1008..0x100b} \\
\code{a[3]} & \code{0x100c} & \code{0x100c..0x100f} \\
\bottomrule
\end{tabular}
\end{center}

\subsubsection{数组名退化为指针}

在很多表达式中，数组名会退化为指向首元素的指针：

\begin{lstlisting}[language=C++]
int a[4];
int* p = a;      // p == &a[0]
\end{lstlisting}

但数组本身不是指针。下面三个 \code{sizeof} 结果不同：

\begin{lstlisting}[language=C++]
int a[4];

sizeof(a);       // 4 * sizeof(int)
sizeof(&a[0]);   // pointer size
sizeof(int*);    // pointer size
\end{lstlisting}

这一区别对性能代码很重要。模板、span、数组引用、指针和长度参数会影响编译器是否知道长度、是否能展开、是否能做边界相关优化。

\subsubsection{越界不是“多走几个字节”这么简单}

从机器地址角度看，\code{a[4]} 似乎只是访问 \code{a[0]} 后面第 16 个字节。但从 C++ 语言角度，访问数组范围外对象是未定义行为。

\begin{lstlisting}[language=C++]
int a[4] = {1, 2, 3, 4};
int x = a[4];    // undefined behavior
\end{lstlisting}

它可能读到某个栈上的值，可能触发崩溃，也可能被编译器假设为永不发生并重写代码。性能工程中不要用越界行为做“技巧”。

\subsection{四步推导法：从 C++ 表达到机器访问}

前面已经讲了指针、数组和类型，但真正读汇编时，你需要一个固定动作。本节给出一个可以反复使用的四步推导法。以后看到任何内存访问表达式，都按这个顺序走：

\begin{enumerate}
  \item \textbf{确定表达式的对象类型和访问宽度。} 它最终读写的是 \code{int}、\code{double}、\code{float}，还是某个结构体字段？访问宽度是 1、2、4、8、16、32 还是 64 字节？
  \item \textbf{确定基地址。} 这个访问从哪个指针、数组首元素、结构体对象或栈帧位置出发？
  \item \textbf{计算字节偏移。} 把下标、字段 offset、数组步长、结构体大小都换成字节单位。
  \item \textbf{映射到 x86-64 寻址。} 尽量写成 \code{[base + index * scale + displacement]}，同时标出 load/store 宽度。
\end{enumerate}

这四步看似简单，但它可以把 C++ 源码、对象布局和汇编连接成一条线。

\subsubsection{例 1：\code{int a[5]; a[3]}}

假设 \code{a[0]} 的地址是 \code{0x1000}，并且 \code{sizeof(int) == 4}。

\begin{center}
\begin{tabular}{p{0.18\linewidth}p{0.68\linewidth}}
\toprule
步骤 & 推导 \\
\midrule
对象类型 & \code{a[3]} 是 \code{int} 左值，访问宽度 4 字节。 \\
基地址 & 数组首元素地址 \code{addr(a[0]) = 0x1000}。 \\
字节偏移 & 下标 3，元素大小 4，所以 offset = \code{3 * 4 = 12 = 0x0c}。 \\
最终地址 & \code{0x1000 + 0x0c = 0x100c}，覆盖字节 \code{0x100c..0x100f}。 \\
\bottomrule
\end{tabular}
\end{center}

如果 \code{a} 作为第一个参数传入函数，\code{i} 作为第二个参数，汇编可能是：

\begin{lstlisting}
mov eax, dword ptr [rdi + rsi*4]
\end{lstlisting}

这里 \register{rdi} 是基地址，\register{rsi} 是元素下标，scale 4 来自 \code{sizeof(int)}，\code{dword ptr} 表示访问 4 字节。

\subsubsection{例 2：\code{double a[5]; a[3]}}

同样是下标 3，但类型换成 \code{double} 后，字节偏移不同。先把条件单独写清楚：

\begin{lstlisting}
addr(a[0])       = 0x1000
sizeof(double)   = 8
\end{lstlisting}

因此：

\begin{lstlisting}
offset = 3 * sizeof(double) = 3 * 8 = 24 = 0x18
addr(a[3]) = 0x1000 + 0x18 = 0x1018
\end{lstlisting}

对应汇编可能是：

\begin{lstlisting}
movsd xmm0, qword ptr [rdi + rsi*8]
\end{lstlisting}

这次 scale 是 8，load width 是 \code{qword}，返回值进入 \register{xmm0}。同一个下标，因为类型不同，机器地址计算不同；同一个地址，因为解释规则不同，得到的值也可能不同。

\subsubsection{例 3：\code{p[i].z} 如何拆成两层偏移}

看一个对 AI/HPC 很常见的结构：

\begin{lstlisting}[language=C++]
struct Vec4 {
    float x;
    float y;
    float z;
    float w;
};

float get_z(Vec4 const* p, std::size_t i)
{
    return p[i].z;
}
\end{lstlisting}

在常见 ABI 下，四个 \code{float} 连续排布：

\begin{center}
\begin{tabular}{llll}
\toprule
字段 & offset & 大小 & 说明 \\
\midrule
\code{x} & 0  & 4 & 第一个 \code{float} \\
\code{y} & 4  & 4 & 第二个 \code{float} \\
\code{z} & 8  & 4 & 第三个 \code{float} \\
\code{w} & 12 & 4 & 第四个 \code{float} \\
\bottomrule
\end{tabular}
\end{center}

所以：

\begin{lstlisting}
addr(p[i].z) = addr(p[0]) + i * sizeof(Vec4) + offset(Vec4, z)
             = base       + i * 16           + 8
\end{lstlisting}

这个地址公式不能直接写成单条 x86-64 内存操作数，因为 x86-64 的 scale 只有 1、2、4、8。真实编译器可能先把 \code{i * 16} 算到寄存器里，再做 load：

\begin{lstlisting}
shl   rsi, 4
movss xmm0, dword ptr [rdi + rsi + 8]
\end{lstlisting}

这里第一条指令把元素下标变成字节偏移。对于 \code{i * 16} 这种 2 的幂乘法，左移通常最直观。对于 12、24、40 这类非 2 的幂步长，后面会专门用 \instruction{lea} 解释。

不要死记某一种汇编形态。关键是看出地址公式里有三个来源：数组基址、元素步长、字段 offset。

\subsubsection{例 4：二维数组下标不是两个 load}

考虑真正的 C++ 二维数组：

\begin{lstlisting}[language=C++]
int a[3][5];
int x = a[i][j];
\end{lstlisting}

\code{a} 是 3 行 5 列的连续数组。按行主序存放，\code{a[i][j]} 的线性元素下标是：

\begin{lstlisting}
i * 5 + j
\end{lstlisting}

字节地址是：

\begin{lstlisting}
addr(a[i][j]) = addr(a[0][0]) + (i * 5 + j) * sizeof(int)
\end{lstlisting}

这不是“先找到第 i 行指针，再从那个指针找第 j 个元素”。真正的 \code{int a[3][5]} 是一整块连续内存，行长度 5 是类型的一部分。只有 \code{int**} 或“指针数组”才是另一种布局。

\begin{warningbox}
\code{int a[3][5]}、\code{int (*p)[5]}、\code{int** p} 是三种不同模型。前两者知道每行 5 个 \code{int}，可以用线性地址公式；\code{int**} 是指向指针的指针，通常需要先读取行指针，再访问该行的数据。它们的性能、局部性和汇编形态都可能不同。
\end{warningbox}

\subsubsection{例 5：把 \code{a[i] += b[i]} 展开成 load-compute-store}

一个赋值表达式往往不止一次内存访问：

\begin{lstlisting}[language=C++]
void add_into(float* a, float const* b, std::size_t n)
{
    for (std::size_t i = 0; i < n; ++i) {
        a[i] += b[i];
    }
}
\end{lstlisting}

循环体在标量层面可以拆成：

\begin{lstlisting}
tmp_a = load 4 bytes from addr(a[0]) + i * 4
tmp_b = load 4 bytes from addr(b[0]) + i * 4
tmp   = tmp_a + tmp_b
store 4 bytes to   addr(a[0]) + i * 4
\end{lstlisting}

这对性能很关键。你不能只数“一次加法”。真实循环每个元素至少有两个 load、一个 store、一个浮点加法、一个循环控制，以及可能的别名检查、边界处理和向量化 remainder。后面做 roofline 模型时，算术强度就是从这种拆解开始的。

\begin{exercisebox}
\textbf{随堂检查 3.1：按四步推导法写答案}

对下面表达式分别写出对象类型、基地址、字节偏移公式和可能的 x86-64 内存操作形式：

\begin{enumerate}
  \item \code{std::uint64_t* p; p[i]}
  \item \code{float* x; x[i + 1]}
  \item \code{struct Pair \{ int a; int b; \}; Pair* p; p[i].b}
  \item \code{double a[4][8]; a[i][j]}
  \item \code{struct Node \{ int key; Node* next; \}; Node* p; p->next}
\end{enumerate}

要求：不要只写 C++ 解释，必须写成字节地址公式。
\end{exercisebox}

\subsection{x86-64 寻址模式：数组访问如何进入指令}

x86-64 内存操作数常见形式是：

\begin{lstlisting}
[base + index * scale + displacement]
\end{lstlisting}

其中：

\begin{itemize}
  \item \code{base} 是基址寄存器。
  \item \code{index} 是索引寄存器。
  \item \code{scale} 通常只能是 1、2、4、8。
  \item \code{displacement} 是编码在指令中的常量偏移。
\end{itemize}

这非常适合表达 C/C++ 的数组和结构体字段访问。

\subsubsection{访问 \code{int a[i]}}

\begin{lstlisting}
mov eax, dword ptr [rdi + rsi*4]
\end{lstlisting}

含义：

\begin{lstlisting}
load 4 bytes from address rdi + rsi * 4
\end{lstlisting}

如果执行到这条指令时寄存器是：

\begin{center}
\begin{tabular}{lll}
\toprule
寄存器 & 值 & 含义 \\
\midrule
\register{rdi} & \code{0x1000} & \code{a[0]} 的地址 \\
\register{rsi} & \code{3} & 下标 \code{i} \\
\bottomrule
\end{tabular}
\end{center}

那么有效地址是：

\begin{lstlisting}
effective_address = 0x1000 + 3 * 4
                  = 0x100c
\end{lstlisting}

\instruction{mov eax, dword ptr [0x100c]} 的语义是从 \code{0x100c} 开始读取 4 字节，写入 \register{eax}。在 x86-64 中，对 \register{eax} 的写入还会清零 \register{rax} 的高 32 位。这一点读汇编时经常重要：\code{int} 返回值写 \register{eax}，但调用者看到的完整 \register{rax} 高位已经被清零。

\subsubsection{访问结构体字段}

如果有：

\begin{lstlisting}[language=C++]
struct S {
    int a;
    double b;
    int c;
};

int get_c(S const* p)
{
    return p->c;
}
\end{lstlisting}

汇编可能是：

\begin{lstlisting}
mov eax, dword ptr [rdi + 16]
ret
\end{lstlisting}

\code{16} 就是字段 \code{c} 相对结构体起始地址的 offset。编译器根据类型布局规则算出它，然后把常量写进指令。

这个例子可以完整展开。常见 x86-64 ABI 下，\code{S} 的布局通常是：

\begin{center}
\begin{tabular}{llll}
\toprule
字节范围 & 内容 & 大小 & 原因 \\
\midrule
0..3 & \code{a} & 4 & \code{int} 字段 \\
4..7 & padding & 4 & 让 \code{b} 满足 8 字节对齐 \\
8..15 & \code{b} & 8 & \code{double} 字段 \\
16..19 & \code{c} & 4 & \code{int} 字段 \\
20..23 & tail padding & 4 & 让 \code{sizeof(S)} 是 8 的倍数 \\
\bottomrule
\end{tabular}
\end{center}

假设 \register{rdi} 保存 \code{p == 0x2000}，那么：

\begin{lstlisting}
addr(p->c) = addr(*p) + offset(S, c)
           = 0x2000   + 16
           = 0x2010
\end{lstlisting}

因此：

\begin{lstlisting}
mov eax, dword ptr [rdi + 16]
\end{lstlisting}

就是从 \code{0x2010..0x2013} 读取 4 字节。这里没有 index 和 scale，因为没有数组下标；只有结构体基地址和字段偏移。

\subsubsection{结构体数组字段访问：scale 和 displacement 同时出现}

更接近真实 kernel 的情况是结构体数组：

\begin{lstlisting}[language=C++]
struct Pixel {
    std::uint8_t r;
    std::uint8_t g;
    std::uint8_t b;
    std::uint8_t a;
};

std::uint8_t load_b(Pixel const* p, std::size_t i)
{
    return p[i].b;
}
\end{lstlisting}

\code{sizeof(Pixel) == 4}，字段 \code{b} 的 offset 是 2，所以地址公式是：

\begin{lstlisting}
addr(p[i].b) = addr(p[0]) + i * sizeof(Pixel) + offset(Pixel, b)
             = base       + i * 4             + 2
\end{lstlisting}

可能的汇编是：

\begin{lstlisting}
movzx eax, byte ptr [rdi + rsi*4 + 2]
ret
\end{lstlisting}

这条指令有几个细节：

\begin{itemize}
  \item \code{byte ptr} 表示从内存读取 1 字节。
  \item \instruction{movzx} 表示 zero extend，读取 8 位后零扩展到 32 位 \register{eax}。
  \item \code{rsi*4} 来自 \code{sizeof(Pixel)}。
  \item \code{+ 2} 来自字段 \code{b} 的 offset。
\end{itemize}

假设：

\begin{center}
\begin{tabular}{lll}
\toprule
寄存器 & 值 & 含义 \\
\midrule
\register{rdi} & \code{0x3000} & \code{p[0]} 的地址 \\
\register{rsi} & \code{5} & 下标 \code{i} \\
\bottomrule
\end{tabular}
\end{center}

则：

\begin{lstlisting}
effective_address = 0x3000 + 5 * 4 + 2
                  = 0x3016
\end{lstlisting}

CPU 会读取 \code{0x3016} 这个地址上的 1 字节，把它零扩展到 32 位返回值。这个案例比 \code{int a[i]} 更接近图像处理、量化张量、packed struct、RGBA buffer 这类真实工作负载。

\subsubsection{当 scale 不够用时：\code{sizeof(T)} 大于 8 怎么办}

x86-64 内存操作数的 scale 只有 1、2、4、8。如果结构体大小是 16、24、32，编译器必须用额外指令先算出字节偏移。

例如：

\begin{lstlisting}[language=C++]
struct Vec4 {
    float x;
    float y;
    float z;
    float w;
};

float load_z(Vec4 const* p, std::size_t i)
{
    return p[i].z;
}
\end{lstlisting}

地址公式是：

\begin{lstlisting}
addr(p[i].z) = base + i * 16 + 8
\end{lstlisting}

一种可能汇编是：

\begin{lstlisting}
shl   rsi, 4
movss xmm0, dword ptr [rdi + rsi + 8]
ret
\end{lstlisting}

如果寄存器初值是：

\begin{center}
\begin{tabular}{lll}
\toprule
寄存器 & 值 & 含义 \\
\midrule
\register{rdi} & \code{0x4000} & \code{p[0]} 的地址 \\
\register{rsi} & \code{3} & 下标 \code{i} \\
\bottomrule
\end{tabular}
\end{center}

执行 \instruction{shl rsi, 4} 后，\register{rsi} 变成 \code{48}，也就是 \code{3 * 16}。随后：

\begin{lstlisting}
effective_address = 0x4000 + 48 + 8
                  = 0x4038
\end{lstlisting}

\instruction{movss} 从 \code{0x4038..0x403b} 读取 4 字节到 \register{xmm0} 的低 32 位。读汇编时要注意：有些寄存器一开始表示“元素下标”，经过 \instruction{shl} 或 \instruction{lea} 后变成“字节偏移”。如果不跟踪这种语义变化，就会把地址算错。

\subsubsection{\instruction{lea}：地址计算指令不一定访问内存}

\instruction{lea} 的名字是 load effective address，但它不从内存读取数据，只做地址形式的整数计算。例如：

\begin{lstlisting}
lea rax, [rsi + rsi*2]
\end{lstlisting}

这条指令计算的是：

\begin{lstlisting}
rax = rsi + rsi * 2
    = 3 * rsi
\end{lstlisting}

它没有访问 \code{[rsi + rsi*2]} 这个内存地址。编译器经常用 \instruction{lea} 做乘法常数、数组偏移和指针计算，因为它可以把“加法 + scale + 常量”合到一条指令里。

例如结构体大小为 24 时：

\begin{lstlisting}[language=C++]
struct Item24 {
    double a;
    double b;
    int c;
    int d;
};

int load_d(Item24 const* p, std::size_t i)
{
    return p[i].d;
}
\end{lstlisting}

若 \code{sizeof(Item24) == 24}，\code{offset(Item24, d) == 20}，地址公式是：

\begin{lstlisting}
addr(p[i].d) = base + i * 24 + 20
\end{lstlisting}

编译器可能生成类似：

\begin{lstlisting}
lea eax, [rsi + rsi*2]        ; eax = i * 3
mov eax, dword ptr [rdi + rax*8 + 20]
ret
\end{lstlisting}

这里第二条指令的 scale 是 8，所以总乘数是 \code{3 * 8 = 24}。这种组合在真实汇编里很常见。读它时不要只看最后一条内存操作，要把前面的 \instruction{lea} 一起纳入地址公式。

\begin{keyidea}
C++ 的数组下标、指针解引用、结构体字段访问，在机器层面都落到地址计算。理解地址计算，是读汇编、理解 cache 行为和优化数据布局的共同入口。
\end{keyidea}

\subsection{对象的存储期：栈、堆、静态区只是第一步分类}

C++ 对象可以有不同存储期。存储期决定对象的存储何时获得、何时释放。常见分类如下：

\begin{center}
\begin{tabular}{lll}
\toprule
存储期 & 示例 & 生命周期大致特点 \\
\midrule
自动存储期 & 局部变量 & 进入作用域创建，离开作用域销毁 \\
动态存储期 & \code{new}、\code{malloc} & 程序显式申请和释放 \\
静态存储期 & 全局变量、\code{static} & 程序启动到结束 \\
线程存储期 & \code{thread_local} & 线程开始到线程结束 \\
\bottomrule
\end{tabular}
\end{center}

初学者常说“局部变量在栈上，new 出来的在堆上，全局变量在静态区”。这作为入门直觉可以，但要知道它不是完整语言规则。编译器优化可能把局部变量放在寄存器里，甚至完全消除；动态分配器背后可能使用 \code{mmap}；静态对象还涉及 ELF 段、重定位和运行时初始化。

\subsubsection{栈对象}

局部自动对象常见实现是在栈帧中分配：

\begin{lstlisting}[language=C++]
void f()
{
    int x = 1;
    int y = 2;
}
\end{lstlisting}

未优化时，编译器可能给 \code{x} 和 \code{y} 分配栈槽。优化后，如果它们不需要真实地址，可能只存在于寄存器甚至常量传播结果中。

\subsubsection{堆对象}

动态对象常由分配器管理：

\begin{lstlisting}[language=C++]
auto* p = new int[1024];
delete[] p;
\end{lstlisting}

堆分配有额外成本：分配器元数据、锁或线程缓存、对齐、碎片、page fault、NUMA 位置等。性能关键路径中频繁分配小对象通常很危险。

\subsubsection{静态对象}

全局变量和静态变量通常在可执行文件或共享库的段中描述，并在进程加载时映射：

\begin{lstlisting}[language=C++]
int global_counter = 0;

int f()
{
    static int local_static = 1;
    return ++local_static;
}
\end{lstlisting}

静态对象会牵涉链接、加载、初始化顺序、线程安全局部静态初始化等问题。本书后面讲 ELF、链接和运行时时会展开。

\subsection{对齐：对象地址不是随便放的}

\term{alignment}{对齐} 是对象地址必须满足的约束。若一个类型 \code{T} 的对齐要求是 \code{A}，则 \code{T} 对象地址通常必须是 \code{A} 的倍数。

例如在常见 x86-64 ABI 下：

\begin{center}
\begin{tabular}{lll}
\toprule
类型 & 常见大小 & 常见对齐 \\
\midrule
\code{char} & 1 & 1 \\
\code{int} & 4 & 4 \\
\code{float} & 4 & 4 \\
\code{double} & 8 & 8 \\
\code{void*} & 8 & 8 \\
\bottomrule
\end{tabular}
\end{center}

你可以用 \code{alignof} 查询：

\begin{lstlisting}[language=C++]
#include <iostream>

int main()
{
    std::cout << alignof(char) << '\n';
    std::cout << alignof(int) << '\n';
    std::cout << alignof(double) << '\n';
    std::cout << alignof(void*) << '\n';
}
\end{lstlisting}

\subsubsection{为什么需要对齐}

对齐来自硬件和 ABI 的共同需要。硬件访问自然对齐的数据通常更简单、更快。有些架构上未对齐访问会触发异常；x86-64 通常允许许多未对齐访问，但可能带来性能代价，尤其是跨 cache line、跨页或用于某些 SIMD 指令时。

例如，一个 8 字节 load 如果地址是 8 字节对齐，且不跨 cache line，硬件处理更直接。如果它从 cache line 最后 4 字节开始，就会跨两个 cache line，可能需要两个 cache line 都可用。

\begin{deepdive}
“x86 支持未对齐访问”不等于“对齐不重要”。标量普通 load/store 可能只是轻微变慢，但跨 cache line、跨 page、原子操作、SIMD 对齐指令、false sharing 场景下，对齐会变成明显性能因素。
\end{deepdive}

\subsection{结构体布局：字段、padding 和 ABI}

结构体不是简单把字段大小相加。编译器需要为每个字段选择 offset，并在必要时插入 \term{padding}{填充字节}，以满足字段和整个结构体的对齐要求。

\subsubsection{先掌握一个机械算法}

学习结构体布局时，不要一上来凭感觉。对普通标准布局结构体，可以先用下面这个机械算法训练。真实 ABI 有更多细节，例如继承、虚函数、空基类优化、位域、向量类型、显式 \code{alignas}、packing pragma，但本算法足够支撑本章和大量性能分析。

\begin{enumerate}
  \item 令当前 offset 为 0。
  \item 从第一个字段开始，按声明顺序处理。
  \item 对当前字段，先查出字段大小 \code{sizeof(field)} 和对齐 \code{alignof(field)}。
  \item 如果当前 offset 不是字段对齐的倍数，就插入 padding，把 offset 向上补齐到下一个满足对齐的位置。
  \item 当前 offset 就是该字段的 offset。
  \item offset 增加字段大小。
  \item 所有字段处理完后，结构体整体对齐通常是所有字段对齐要求的最大值。
  \item 最终 \code{sizeof(struct)} 必须补齐到结构体整体对齐的倍数，因此末尾可能有 tail padding。
\end{enumerate}

其中“向上补齐”可以写成一个函数：

\begin{center}
\begin{minipage}{0.72\linewidth}
\begin{verbatim}
align_up(x, A) = smallest y such that y >= x and y % A == 0
\end{verbatim}
\end{minipage}
\end{center}

如果 \code{A} 是 2 的幂，机器代码和系统代码里经常写成：

\begin{lstlisting}[language=C++]
std::size_t align_up(std::size_t x, std::size_t a)
{
    return (x + a - 1) & ~(a - 1);
}
\end{lstlisting}

例如 \code{align_up(5, 4) == 8}，\code{align_up(16, 8) == 16}。第一个例子需要补 3 字节，第二个例子本来已经对齐，不需要补。

\begin{keyidea}
padding 不是编译器“随便浪费空间”。padding 是为了让每个字段和数组中每个结构体元素都满足对齐约束。结构体大小必须包含尾部 padding，否则 \code{arr[1]} 的起始地址可能不满足结构体自身对齐。
\end{keyidea}

看第一个例子：

\begin{lstlisting}[language=C++]
struct A {
    char c;
    int x;
};
\end{lstlisting}

在常见 x86-64 ABI 下，\code{char} 大小 1、对齐 1；\code{int} 大小 4、对齐 4。若 \code{c} 放在 offset 0，那么 \code{x} 不能放在 offset 1，因为 offset 1 不是 4 的倍数。编译器会插入 3 字节 padding：

\begin{center}
\begin{tabular}{llll}
\toprule
offset & 内容 & 大小 & 说明 \\
\midrule
0 & \code{c} & 1 & 字段 \\
1..3 & padding & 3 & 让下个字段 4 字节对齐 \\
4..7 & \code{x} & 4 & 字段 \\
\bottomrule
\end{tabular}
\end{center}

所以 \code{sizeof(A)} 通常是 8，不是 5。

\subsubsection{完整例题：逐字段模拟布局}

现在用算法推一个稍复杂的例子：

\begin{lstlisting}[language=C++]
struct Record {
    char tag;
    double score;
    int count;
    short flags;
};
\end{lstlisting}

假设常见 x86-64 ABI 下：

\begin{center}
\begin{tabular}{lll}
\toprule
类型 & 大小 & 对齐 \\
\midrule
\code{char} & 1 & 1 \\
\code{short} & 2 & 2 \\
\code{int} & 4 & 4 \\
\code{double} & 8 & 8 \\
\bottomrule
\end{tabular}
\end{center}

逐字段推导如下：

\begin{center}
\begin{tabular}{lllll}
\toprule
字段 & 进入前 offset & 对齐动作 & 字段 offset & 写入后 offset \\
\midrule
\code{tag} & 0 & 已满足 1 对齐 & 0 & 1 \\
\code{score} & 1 & 补到 8，插入 7 字节 padding & 8 & 16 \\
\code{count} & 16 & 已满足 4 对齐 & 16 & 20 \\
\code{flags} & 20 & 已满足 2 对齐 & 20 & 22 \\
\bottomrule
\end{tabular}
\end{center}

字段处理完后，最大对齐是 8，所以结构体大小必须补到 8 的倍数。当前 offset 是 22，\code{align_up(22, 8) == 24}，因此尾部 padding 是 2 字节：

\begin{center}
\begin{tabular}{llll}
\toprule
字节范围 & 内容 & 大小 & 说明 \\
\midrule
0 & \code{tag} & 1 & 字段 \\
1..7 & padding & 7 & 让 \code{score} 8 字节对齐 \\
8..15 & \code{score} & 8 & 字段 \\
16..19 & \code{count} & 4 & 字段 \\
20..21 & \code{flags} & 2 & 字段 \\
22..23 & tail padding & 2 & 让 \code{sizeof(Record)} 是 8 的倍数 \\
\bottomrule
\end{tabular}
\end{center}

如果有数组：

\begin{lstlisting}[language=C++]
Record r[2];
\end{lstlisting}

那么：

\begin{center}
\begin{minipage}{0.72\linewidth}
\begin{verbatim}
addr(r[1]) = addr(r[0]) + sizeof(Record)
           = addr(r[0]) + 24
\end{verbatim}
\end{minipage}
\end{center}

\code{24} 是 8 的倍数，因此 \code{r[1].score} 仍然可以落在 8 字节对齐地址上。如果结构体大小只是字段结束时的 22，数组第二个元素就会破坏字段对齐，这就是尾部 padding 存在的根本原因。

\subsubsection{把字段重排也算法化}

如果只是追求更小的结构体大小，一个常见启发式是把对齐要求大的字段放前面：

\begin{lstlisting}[language=C++]
struct RecordPackedOrder {
    double score;
    int count;
    short flags;
    char tag;
};
\end{lstlisting}

逐字段推导：

\begin{center}
\begin{tabular}{lllll}
\toprule
字段 & 进入前 offset & 对齐动作 & 字段 offset & 写入后 offset \\
\midrule
\code{score} & 0 & 已满足 8 对齐 & 0 & 8 \\
\code{count} & 8 & 已满足 4 对齐 & 8 & 12 \\
\code{flags} & 12 & 已满足 2 对齐 & 12 & 14 \\
\code{tag} & 14 & 已满足 1 对齐 & 14 & 15 \\
\bottomrule
\end{tabular}
\end{center}

最大对齐仍然是 8，最终大小补到 16。和原来的 24 相比，减少了 8 字节。如果这种对象有一千万个，节省的是约 80 MB，不是小数。

但这仍然不是“所有结构体都按大到小排序”的命令。性能结构体设计要同时考虑三类目标：

\begin{itemize}
  \item \textbf{空间密度：} 减少 padding，让更多对象进入 cache。
  \item \textbf{访问局部性：} 把同一个 hot loop 经常一起使用的字段放近，把冷字段挪远。
  \item \textbf{并发隔离：} 多线程频繁写的字段有时要故意分开，避免 false sharing。
\end{itemize}

字段顺序是数据布局设计的一部分，不只是“省几个字节”。对于 AI runtime 的 tensor metadata、operator descriptor、thread task、work queue node，字段布局会影响 cache 占用、预取效果、false sharing 和 ABI 稳定性。

\subsubsection{尾部 padding}

再看：

\begin{lstlisting}[language=C++]
struct B {
    int x;
    char c;
};
\end{lstlisting}

\code{x} 放在 offset 0，占 4 字节；\code{c} 放在 offset 4，占 1 字节。看起来总共 5 字节。但结构体整体对齐是其成员最大对齐，通常是 4。为了让 \code{B arr[2]} 中每个元素都满足对齐，\code{sizeof(B)} 必须是 4 的倍数，于是尾部会有 3 字节 padding：

\begin{center}
\begin{tabular}{llll}
\toprule
offset & 内容 & 大小 & 说明 \\
\midrule
0..3 & \code{x} & 4 & 字段 \\
4 & \code{c} & 1 & 字段 \\
5..7 & tail padding & 3 & 让数组中下个元素对齐 \\
\bottomrule
\end{tabular}
\end{center}

\subsubsection{字段顺序会影响大小}

比较：

\begin{lstlisting}[language=C++]
struct Bad {
    char a;
    double b;
    char c;
    int d;
};

struct Better {
    double b;
    int d;
    char a;
    char c;
};
\end{lstlisting}

\code{Bad} 可能因为多次对齐插入大量 padding；\code{Better} 把大对齐字段放前面，通常更紧凑。

但字段重排不是无脑优化。你还要考虑：

\begin{itemize}
  \item ABI 兼容：公开结构体布局不能随便改。
  \item 序列化和文件格式：布局变化会破坏二进制兼容。
  \item 热字段和冷字段：有时为了 cache，把常用字段聚在一起比最小 \code{sizeof} 更重要。
  \item false sharing：多线程写的字段可能需要故意分离到不同 cache line。
\end{itemize}

\subsection{\code{sizeof}、\code{alignof} 和 \code{offsetof}}

研究布局时，三个工具非常重要：

\begin{lstlisting}[language=C++]
#include <cstddef>
#include <iostream>

struct S {
    char a;
    double b;
    int c;
};

int main()
{
    std::cout << "sizeof(S)  = " << sizeof(S) << '\n';
    std::cout << "alignof(S) = " << alignof(S) << '\n';
    std::cout << "offset a   = " << offsetof(S, a) << '\n';
    std::cout << "offset b   = " << offsetof(S, b) << '\n';
    std::cout << "offset c   = " << offsetof(S, c) << '\n';
}
\end{lstlisting}

\code{offsetof} 对标准布局类型可靠。学习阶段你应该先手工预测，再用程序验证。不要只看输出，因为真正的训练是形成布局推导能力。

\subsection{padding 字节不一定有稳定值}

结构体 padding 是对象表示的一部分，但不一定参与对象的值。它可能未初始化，也可能包含旧栈数据。比较结构体对象时，不能简单用 \code{memcmp} 替代字段级比较，除非你非常清楚类型性质和初始化方式。

例如：

\begin{lstlisting}[language=C++]
struct A {
    char c;
    int x;
};

bool equal_bad(A const& lhs, A const& rhs)
{
    return std::memcmp(&lhs, &rhs, sizeof(A)) == 0;
}
\end{lstlisting}

即使两个对象的 \code{c} 和 \code{x} 相等，padding 字节也可能不同，\code{memcmp} 可能返回不相等。这个问题在哈希、序列化、网络传输和 SIMD 批量比较中很常见。

\begin{warningbox}
结构体的“值”和结构体的“对象字节”不是总能一一对应。padding、未初始化字节、浮点 NaN、不同但等价的表示，都可能让字节级比较和语义级比较不一致。
\end{warningbox}

\subsection{安全地观察和搬运字节}

\subsubsection{用 \code{unsigned char} 或 \code{std::byte} 观察对象表示}

观察对象字节时，可以使用 \code{unsigned char const*}：

\begin{lstlisting}[language=C++]
template <class T>
void dump_bytes(T const& value)
{
    auto const* p = reinterpret_cast<unsigned char const*>(&value);
    for (std::size_t i = 0; i < sizeof(T); ++i) {
        std::cout << std::hex << static_cast<unsigned>(p[i]) << ' ';
    }
    std::cout << '\n';
}
\end{lstlisting}

\subsubsection{用 \code{std::memcpy} 或 \code{std::bit_cast} 搬运位模式}

如果你想在两种类型之间复制位模式，不要随便用错误类型指针解引用。C++20 提供 \code{std::bit_cast}：

\begin{lstlisting}[language=C++]
#include <bit>
#include <cstdint>

float f = 1.0f;
std::uint32_t bits = std::bit_cast<std::uint32_t>(f);
\end{lstlisting}

也可以用 \code{std::memcpy}：

\begin{lstlisting}[language=C++]
std::uint32_t bits;
std::memcpy(&bits, &f, sizeof(bits));
\end{lstlisting}

这两种方式都表达“复制对象表示”，比用不匹配类型的指针解引用更可靠。

\subsection{别名规则：编译器如何知道两个指针会不会指同一块内存}

别名问题是从 C++ 到高性能优化的关键门槛。看下面代码：

\begin{lstlisting}[language=C++]
void f(int* a, int* b)
{
    *a = 1;
    *b = 2;
    *a = *a + 3;
}
\end{lstlisting}

如果 \code{a} 和 \code{b} 指向不同对象，最后 \code{*a == 4}。如果它们指向同一个 \code{int}，执行过程是：

\begin{enumerate}
  \item 执行 \code{*a = 1}，该对象变成 1。
  \item 执行 \code{*b = 2}。因为 \code{a} 和 \code{b} 指向同一对象，所以 \code{*a} 也变成 2。
  \item 执行 \code{*a = *a + 3}，最终该对象变成 5。
\end{enumerate}

编译器如果不能证明 \code{a} 和 \code{b} 不别名，就必须保守处理。保守处理会影响寄存器缓存、load/store 重排、循环向量化和 common subexpression elimination。

\subsubsection{不同类型通常提供别名信息}

看另一个例子：

\begin{lstlisting}[language=C++]
void g(int* a, double* b)
{
    *a = 1;
    *b = 2.0;
    *a = *a + 3;
}
\end{lstlisting}

在正常 C++ 别名规则下，\code{int*} 和 \code{double*} 通常不能指向同一个活着的对象。编译器可以利用这个事实，把 \code{*a} 保存在寄存器中，而不必担心 \code{*b = 2.0} 改了同一个 \code{int}。

这就是为什么错误的类型双关不仅可能不正确，还可能导致编译器生成你无法预期的代码。

\begin{deepdive}
C 的 \code{restrict}、C++ 编译器扩展中的 \code{__restrict__}、LLVM IR 的 alias metadata、Fortran 的数组别名语义，都会显著影响高性能循环能否向量化。后面讲 auto-vectorization 和 AI kernel 时会深入。
\end{deepdive}

\subsection{常见混淆：底层学习最容易卡住的地方}

本章内容多，但很多错误可以归结为几组概念混淆。这里集中整理一次。你后面看汇编、写 SIMD、调 benchmark 时，只要结果不符合预期，就应该回到这些问题检查。

\subsubsection{混淆 1：指针变量和被指对象不是同一个东西}

看下面代码：

\begin{lstlisting}[language=C++]
int x = 42;
int* p = &x;
\end{lstlisting}

\code{x} 是一个 \code{int} 对象，通常占 4 字节。\code{p} 也是一个对象，它的类型是 \code{int*}，在 x86-64 上通常占 8 字节。\code{p} 的值是 \code{x} 的地址，但 \code{p} 自己也需要存储空间。

如果假设：

\begin{lstlisting}
addr(x) = 0x1000
addr(p) = 0x2000
value(p) = 0x1000
\end{lstlisting}

那么：

\begin{itemize}
  \item 访问 \code{x}，读取的是 \code{0x1000..0x1003}。
  \item 访问 \code{p} 这个指针变量，读取的是 \code{0x2000..0x2007}。
  \item 访问 \code{*p}，先从 \code{p} 的值拿到 \code{0x1000}，再读取 \code{0x1000..0x1003}。
\end{itemize}

汇编里也会体现这个差异。如果 \register{rdi} 直接保存 \code{p} 的值，也就是 \code{x} 的地址，那么：

\begin{lstlisting}
mov eax, dword ptr [rdi]
\end{lstlisting}

读取的是 \code{*p}。但如果 \register{rdi} 保存的是指针变量 \code{p} 自己的地址，也就是 \code{&p}，那么通常需要两次 load：

\begin{lstlisting}
mov rax, qword ptr [rdi]      ; load p itself
mov eax, dword ptr [rax]      ; load *p
\end{lstlisting}

第一条读取指针值，第二条读取被指对象。链表、树、稀疏矩阵和 \code{int**} 的性能问题，很多都来自这种“先读地址，再按地址读数据”的间接访问。

\subsubsection{混淆 2：元素偏移和字节偏移不是同一个单位}

下面两行很像，但单位不同：

\begin{lstlisting}[language=C++]
int* q1 = p + 3;
auto q2 = reinterpret_cast<char*>(p) + 3;
\end{lstlisting}

如果 \code{p == 0x1000}，则：

\begin{lstlisting}
q1 = 0x1000 + 3 * sizeof(int) = 0x100c
q2 = 0x1000 + 3                = 0x1003
\end{lstlisting}

\code{p + 3} 的 3 是元素个数；\code{reinterpret_cast<char*>(p) + 3} 的 3 是字节数。读汇编时，所有地址最终都要落到字节偏移；读 C++ 时，指针加法先按类型缩放。

\subsubsection{混淆 3：数组对象、数组首元素指针、指向整个数组的指针}

看代码：

\begin{lstlisting}[language=C++]
int a[4];
int* p0 = a;
int (*p1)[4] = &a;
\end{lstlisting}

它们的数值地址可能相同，但类型不同：

\begin{center}
\begin{tabular}{llll}
\toprule
表达式 & 类型 & 数值上常见地址 & \code{+1} 的含义 \\
\midrule
\code{a} & 表达式中常退化为 \code{int*} & \code{addr(a[0])} & 下一个 \code{int} \\
\code{&a[0]} & \code{int*} & \code{addr(a[0])} & 下一个 \code{int} \\
\code{&a} & \code{int (*)[4]} & \code{addr(a)} & 下一个 4 元素数组 \\
\bottomrule
\end{tabular}
\end{center}

因此：

\begin{lstlisting}
addr((&a[0]) + 1) = base + 4
addr((&a)    + 1) = base + 16
\end{lstlisting}

这类区别在多维数组、固定 shape 张量、小矩阵 tile、编译期已知长度的 kernel 中非常重要。编译器知道更多类型信息时，往往能生成更直接的地址计算和更强的优化。

\subsubsection{混淆 4：硬件能执行不代表 C++ 定义良好}

很多初学者会用“我运行没崩”来判断底层代码对不对。这在系统编程和性能优化里非常危险。看下面例子：

\begin{lstlisting}[language=C++]
int a[4] = {1, 2, 3, 4};
int x = a[4];
\end{lstlisting}

硬件层面，这可能只是从 \code{a[0]} 后面 16 字节的位置发起一次 4 字节 load。但 C++ 层面，\code{a[4]} 越界，行为未定义。编译器可以假设这种情况不会发生，并据此重排、删除或改写代码。

再看类型双关：

\begin{lstlisting}[language=C++]
int x = 0;
double* p = reinterpret_cast<double*>(&x);
double y = *p;
\end{lstlisting}

硬件也许能从 \code{&x} 开始读 8 字节，但 C++ 对象模型、对齐和别名规则都可能已经被破坏。正确观察位模式时，应优先使用：

\begin{lstlisting}[language=C++]
std::bit_cast<T>(value)
std::memcpy(&dst, &src, sizeof(dst))
unsigned char const* bytes = reinterpret_cast<unsigned char const*>(&obj)
\end{lstlisting}

底层工程的目标不是“骗过编译器和硬件”，而是在语言规则允许的范围内，把编译器、ISA 和微架构都用到极限。

\subsubsection{混淆 5：汇编里的地址公式不等于源码里的唯一写法}

同一个源码可能生成多种汇编，同一个汇编片段也可能来自多种源码。比如：

\begin{lstlisting}
mov eax, dword ptr [rdi + rsi*4 + 8]
\end{lstlisting}

它可能来自：

\begin{lstlisting}[language=C++]
return p[i + 2];
\end{lstlisting}

也可能来自：

\begin{lstlisting}[language=C++]
struct S {
    int a;
    int b;
    int c;
};

return p[i].c;
\end{lstlisting}

区别要靠上下文判断：\register{rdi} 是不是数组首地址？\register{rsi} 是不是原始下标？结构体大小是否为 12？前面是否有 \instruction{lea} 改变了 index？函数参数 ABI、调试信息、周围 load/store 和循环结构都要一起看。

\begin{mentalmodel}
读汇编不是把每条指令机械翻译回一行 C++。更可靠的方法是先恢复数据流和地址公式，再结合类型、ABI、循环结构和访问宽度推断源码意图。
\end{mentalmodel}

\subsection{对象生命周期：有地址不等于有对象}

底层程序员容易犯的另一个错误是把“有一块字节存储”当成“有一个对象”。C++ 区分存储和对象生命周期。

例如，分配一块原始存储：

\begin{lstlisting}[language=C++]
alignas(int) unsigned char buffer[sizeof(int)];
\end{lstlisting}

这块字节存在，但其中还没有一个 \code{int} 对象。要在其中创建对象，需要 placement new 或符合标准的对象创建方式：

\begin{lstlisting}[language=C++]
int* p = new (buffer) int(42);
\end{lstlisting}

在高性能代码中，自定义 allocator、memory pool、arena、SIMD buffer 和序列化反序列化都会碰到对象生命周期问题。初学阶段不要求你立刻掌握所有细节，但必须建立警觉：地址、字节、对象不是同一个概念。

\subsection{数据布局如何进入性能：cache line 预告}

CPU 不是每次只从内存拿你请求的那几个字节。缓存层级通常以 cache line 为单位搬运数据。现代 x86-64 常见 cache line 大小是 64 字节。

如果你顺序遍历 \code{int} 数组：

\begin{lstlisting}[language=C++]
for (std::size_t i = 0; i < n; ++i) {
    s += a[i];
}
\end{lstlisting}

每个 \code{int} 4 字节，一个 64 字节 cache line 可以容纳 16 个 \code{int}。当 CPU 加载 \code{a[i]} 所在 cache line 后，后面的 \code{a[i+1]} 到 \code{a[i+15]} 很可能已经在 cache 中。这叫空间局部性。

如果你随机访问：

\begin{lstlisting}[language=C++]
for (std::size_t i = 0; i < n; ++i) {
    s += a[index[i]];
}
\end{lstlisting}

每次访问可能落在不同 cache line。即使只需要 4 字节，也可能为它搬来 64 字节，造成带宽浪费和高延迟。

\subsubsection{步长访问}

步长访问是理解 cache 的第一个实验：

\begin{lstlisting}[language=C++]
for (std::size_t i = 0; i < n; i += stride) {
    s += a[i];
}
\end{lstlisting}

当 \code{stride == 1} 时，空间局部性最好。当 \code{stride} 很大时，每次访问都可能使用一个新 cache line。访问元素数看似减少了，但每个访问的 cache line 利用率也下降了。

\subsection{AoS 与 SoA：布局决定 SIMD 和 cache 效率}

看一个粒子数组：

\begin{lstlisting}[language=C++]
struct Particle {
    float x;
    float y;
    float z;
    float mass;
};

Particle* particles;
\end{lstlisting}

这是 \term{AoS}{Array of Structs}：数组中的每个元素是一个结构体。如果只计算所有 \code{x} 的和：

\begin{lstlisting}[language=C++]
float sx = 0.0f;
for (std::size_t i = 0; i < n; ++i) {
    sx += particles[i].x;
}
\end{lstlisting}

内存中 \code{x} 被 \code{y}、\code{z}、\code{mass} 间隔开。每次 cache line 搬入很多不需要的字段。

另一种布局是 \term{SoA}{Struct of Arrays}：

\begin{lstlisting}[language=C++]
struct Particles {
    float* x;
    float* y;
    float* z;
    float* mass;
};
\end{lstlisting}

此时所有 \code{x} 连续存放。对只读 \code{x} 的 kernel，SoA 更利于 cache 和 SIMD。

但如果每次都同时使用 \code{x,y,z,mass}，AoS 可能更直观，也可能有不错局部性。真正的布局选择必须围绕访问模式，而不是口号。

\begin{keyidea}
数据布局不是“代码风格”，而是性能模型的一部分。一个 kernel 的访问模式决定它需要什么布局；布局决定 cache line 利用率、SIMD load/store 形态、预取效果和多线程 false sharing 风险。
\end{keyidea}

\subsection{从布局回到汇编：offset 和 scale 是布局的影子}

读反汇编时，不要只盯着某一条 \instruction{mov}。你要把它前面的地址计算指令也一起看。下面几组模式非常常见。

\subsubsection{模式 1：纯数组，scale 直接等于元素大小}

当你看到：

\begin{lstlisting}
movss xmm0, dword ptr [rdi + rax*4]
\end{lstlisting}

它很可能表示读取一个连续 \code{float} 数组：

\begin{lstlisting}[language=C++]
float x = arr[i];
\end{lstlisting}

解释如下：

\begin{center}
\begin{tabular}{lll}
\toprule
汇编部分 & 含义 & 可能来源 \\
\midrule
\register{rdi} & 基地址 & \code{arr} \\
\register{rax} & 下标 & \code{i} \\
\code{*4} & 元素大小 & \code{sizeof(float)} \\
\code{dword ptr} & 访问 4 字节 & \code{float} load \\
\bottomrule
\end{tabular}
\end{center}

\subsubsection{模式 2：结构体数组字段，先算大步长}

如果结构体大小是 16，x86 内存操作数不能直接写 \code{rax*16}。真实汇编更可能是：

\begin{lstlisting}
shl   rax, 4
movss xmm0, dword ptr [rdi + rax + 4]
\end{lstlisting}

它可能表示：

\begin{lstlisting}[language=C++]
struct Particle {
    float x;
    float y;
    float z;
    float mass;
};

float y = particles[i].y;
\end{lstlisting}

推导过程是：

\begin{lstlisting}
sizeof(Particle)       = 16
offset(Particle, y)    = 4
byte_offset_for_i      = i * 16
addr(particles[i].y)   = base + i * 16 + 4
\end{lstlisting}

所以 \instruction{shl rax, 4} 把元素下标变成字节偏移，\code{+ 4} 是字段 \code{y} 的 offset。

\subsubsection{模式 3：用 \instruction{lea} 拼出非 2 的幂步长}

当结构体大小是 24 时，常见组合是：

\begin{lstlisting}
lea   rcx, [rax + rax*2]
mov   edx, dword ptr [rdi + rcx*8 + 20]
\end{lstlisting}

这可能表示：

\begin{lstlisting}[language=C++]
struct Item24 {
    double a;
    double b;
    int c;
    int d;
};

int d = items[i].d;
\end{lstlisting}

因为：

\begin{lstlisting}
rcx = rax + rax * 2 = i * 3
effective_address = base + (i * 3) * 8 + 20
                  = base + i * 24 + 20
\end{lstlisting}

\code{20} 是字段 \code{d} 的 offset。这个例子训练的是：不要看到最后一条里只有 \code{rcx*8} 就误判结构体大小是 8；\register{rcx} 本身已经被前一条 \instruction{lea} 变成了 \code{i*3}。

\subsubsection{模式 4：二维连续数组，行长进入地址公式}

如果源码是：

\begin{lstlisting}[language=C++]
int get_2d(int const a[3][5], std::size_t i, std::size_t j)
{
    return a[i][j];
}
\end{lstlisting}

一种可能汇编形态是：

\begin{lstlisting}
lea eax, [rsi + rsi*4]
add rax, rdx
mov eax, dword ptr [rdi + rax*4]
ret
\end{lstlisting}

其中：

\begin{lstlisting}
rax = i + i * 4 = i * 5
rax = i * 5 + j
address = base + (i * 5 + j) * 4
\end{lstlisting}

这说明每行 5 个 \code{int} 这个类型信息已经进入机器地址计算。如果是 \code{int**}，则通常会先 load 行指针，再访问行内元素，汇编形态完全不同。

\subsubsection{反汇编阅读清单}

看到一条内存操作时，按下面清单写笔记：

\begin{enumerate}
  \item 这条指令是 load 还是 store？
  \item 访问宽度是多少：\code{byte}、\code{word}、\code{dword}、\code{qword}、\code{xmmword} 还是 \code{ymmword}？
  \item base 寄存器可能是哪一个指针或对象起点？
  \item index 寄存器现在是元素下标，还是已经被 \instruction{shl}/\instruction{lea} 变成字节偏移？
  \item scale 来自元素大小、字段数组大小，还是地址计算的中间步骤？
  \item displacement 是结构体字段 offset、固定下标、栈槽偏移，还是全局对象偏移？
  \item 如果有前置 \instruction{lea}，它把下标乘成了多少？
  \item 这次访问落到一个 cache line 内，还是可能跨 cache line？
\end{enumerate}

\begin{keyidea}
offset 和 scale 是数据布局在汇编里的影子；\instruction{shl} 和 \instruction{lea} 则经常是“影子被折叠之前”的地址计算过程。读 AI 算子汇编时，先恢复地址公式，再讨论 SIMD、cache、带宽和流水线。
\end{keyidea}

\subsection{实验 3.1：打印地址并区分存储期}

\begin{labbox}
创建 \filepath{labs/ch03_data_layout/address_map.cpp}：

\begin{lstlisting}[language=C++]
#include <cstdlib>
#include <iostream>
#include <memory>

int global_x = 1;

int main()
{
    static int static_x = 2;
    int stack_x = 3;
    auto heap_x = std::make_unique<int>(4);

    std::cout << "&global_x = " << &global_x << '\n';
    std::cout << "&static_x = " << &static_x << '\n';
    std::cout << "&stack_x  = " << &stack_x << '\n';
    std::cout << "heap_x    = " << heap_x.get() << '\n';

    int stack_arr[4] = {1, 2, 3, 4};
    for (int i = 0; i < 4; ++i) {
        std::cout << "&stack_arr[" << i << "] = "
                  << &stack_arr[i] << '\n';
    }
}
\end{lstlisting}

编译运行：

\begin{lstlisting}[language=bash]
clang++ -std=c++20 -O0 -g address_map.cpp -o address_map
./address_map
./address_map
\end{lstlisting}

交付物：

\begin{enumerate}
  \item 记录两次运行的地址，说明哪些地址变化明显。
  \item 解释栈数组相邻元素地址差值为什么等于 \code{sizeof(int)}。
  \item 说明这些地址为什么是虚拟地址，不应直接理解为物理内存地址。
\end{enumerate}
\end{labbox}

\subsection{实验 3.2：对象字节、字节序和 \code{std::bit_cast}}

\begin{labbox}
创建 \filepath{labs/ch03_data_layout/dump_bytes.cpp}，实现一个通用 \code{dump_bytes}：

\begin{lstlisting}[language=C++]
#include <bit>
#include <cstdint>
#include <cstring>
#include <iomanip>
#include <iostream>
#include <type_traits>

template <class T>
void dump_bytes(char const* name, T const& value)
{
    static_assert(std::is_trivially_copyable_v<T>);

    auto const* p = reinterpret_cast<unsigned char const*>(&value);
    std::cout << name << " (" << sizeof(T) << " bytes): ";
    for (std::size_t i = 0; i < sizeof(T); ++i) {
        std::cout << std::hex << std::setw(2) << std::setfill('0')
                  << static_cast<unsigned>(p[i]) << ' ';
    }
    std::cout << std::dec << '\n';
}

struct A {
    char c;
    int x;
};

int main()
{
    std::uint32_t u = 0x12345678U;
    float f = 1.0f;
    A a{'Z', 0x12345678};

    dump_bytes("u", u);
    dump_bytes("f", f);
    dump_bytes("a", a);

    auto f_bits = std::bit_cast<std::uint32_t>(f);
    dump_bytes("f_bits", f_bits);
}
\end{lstlisting}

交付物：

\begin{enumerate}
  \item 标出 \code{u} 的小端字节序。
  \item 解释 \code{A} 中哪些字节是字段，哪些可能是 padding。
  \item 比较 \code{f} 和 \code{f_bits} 的输出，说明 \code{bit_cast} 的含义。
  \item 用 \code{-O0} 和 \code{-O3} 都运行一次，观察输出是否变化。
\end{enumerate}
\end{labbox}

\subsection{实验 3.3：手工预测结构体布局}

\begin{labbox}
对下面结构体，先手工预测 \code{sizeof}、\code{alignof} 和每个字段 offset，再写程序验证：

\begin{lstlisting}[language=C++]
struct S1 {
    char a;
    int b;
    double c;
    char d;
};

struct S2 {
    double c;
    int b;
    char a;
    char d;
};

struct S3 {
    char tag;
    float x;
    float y;
    float z;
};
\end{lstlisting}

验证程序必须使用：

\begin{lstlisting}[language=C++]
sizeof(T)
alignof(T)
offsetof(T, field)
\end{lstlisting}

交付物：

\begin{enumerate}
  \item 每个结构体的布局表。
  \item 手工预测和实际输出的对比。
  \item 解释字段重排为什么改变或没有改变 \code{sizeof}。
  \item 说明如果这些结构体已经写入二进制文件，为什么不能随便重排字段。
\end{enumerate}
\end{labbox}

\subsection{实验 3.4：指针加法和汇编寻址模式}

\begin{labbox}
创建 \filepath{labs/ch03_data_layout/pointer_scale.cpp}。本实验不是只看一条 load 指令，而是训练你把源码、结构体布局、汇编寻址和具体地址模拟连起来。

\begin{lstlisting}[language=C++]
#include <cstddef>
#include <cstdint>

extern "C" int load_int(int const* p, int i)
{
    return p[i];
}

extern "C" double load_double(double const* p, int i)
{
    return p[i];
}

struct Pair {
    int a;
    int b;
};

extern "C" int load_pair_b(Pair const* p, int i)
{
    return p[i].b;
}

struct Pixel {
    std::uint8_t r;
    std::uint8_t g;
    std::uint8_t b;
    std::uint8_t a;
};

extern "C" std::uint8_t load_pixel_b(Pixel const* p, std::size_t i)
{
    return p[i].b;
}

struct Vec4 {
    float x;
    float y;
    float z;
    float w;
};

extern "C" float load_vec4_z(Vec4 const* p, std::size_t i)
{
    return p[i].z;
}

extern "C" int load_2d(int const (*a)[5], std::size_t i, std::size_t j)
{
    return a[i][j];
}
\end{lstlisting}

生成汇编：

\begin{lstlisting}[language=bash]
clang++ -std=c++20 -O3 -S -masm=intel pointer_scale.cpp -o pointer_scale.s
clang++ -std=c++20 -O3 -c pointer_scale.cpp -o pointer_scale.o
objdump -d -Mintel pointer_scale.o > pointer_scale.objdump.txt
\end{lstlisting}

如果系统上没有 GNU \code{objdump}，可以尝试：

\begin{lstlisting}[language=bash]
llvm-objdump -d --x86-asm-syntax=intel pointer_scale.o > pointer_scale.objdump.txt
\end{lstlisting}

第一部分交付物：为每个函数写一张汇编阅读表。

\begin{center}
\begin{tabular}{p{0.19\linewidth}p{0.18\linewidth}p{0.16\linewidth}p{0.16\linewidth}p{0.18\linewidth}}
\toprule
函数 & load 指令 & base & index/scale & displacement \\
\midrule
\code{load_int} & 你填写 & 你填写 & 你填写 & 你填写 \\
\code{load_double} & 你填写 & 你填写 & 你填写 & 你填写 \\
\code{load_pair_b} & 你填写 & 你填写 & 你填写 & 你填写 \\
\code{load_pixel_b} & 你填写 & 你填写 & 你填写 & 你填写 \\
\code{load_vec4_z} & 你填写 & 你填写 & 你填写 & 你填写 \\
\code{load_2d} & 你填写 & 你填写 & 你填写 & 你填写 \\
\bottomrule
\end{tabular}
\end{center}

第二部分交付物：为每个函数写地址公式。格式必须像下面这样明确：

\begin{lstlisting}
load_pair_b:
    sizeof(Pair) = 8
    offset(Pair, b) = 4
    addr(p[i].b) = base + i * 8 + 4

load_vec4_z:
    sizeof(Vec4) = 16
    offset(Vec4, z) = 8
    addr(p[i].z) = base + i * 16 + 8
\end{lstlisting}

第三部分交付物：选择三个函数做具体地址模拟。至少包含一个 scale 为 4 或 8 的例子、一个结构体字段例子、一个需要 \instruction{shl} 或 \instruction{lea} 的例子。

示例格式：

\begin{lstlisting}
Function: load_pixel_b
Assume:
    rdi = 0x3000
    rsi = 5
Formula:
    addr(p[i].b) = 0x3000 + 5 * 4 + 2
                 = 0x3016
Access:
    load 1 byte from 0x3016
\end{lstlisting}

第四部分交付物：记录目标文件中的机器码字节。报告里必须贴出每个函数的反汇编片段，例如：

\begin{lstlisting}
0000000000000000 <load_int>:
   0:  48 63 f6              movsxd rsi, esi
   3:  8b 04 b7              mov    eax,DWORD PTR [rdi+rsi*4]
   6:  c3                    ret
\end{lstlisting}

如果你的机器生成的指令不同，不要照抄本书示例。你要解释自己的编译器为什么可能选择不同指令，例如用 \instruction{lea}、\instruction{movsxd}、\instruction{shl}、不同寄存器或不同指令顺序。

最终报告必须回答：

\begin{enumerate}
  \item 哪些函数的元素大小可以直接放进 scale？
  \item 哪些函数需要额外地址计算？额外地址计算由哪条指令完成？
  \item 哪些 displacement 来自字段 offset？
  \item \instruction{movzx} 和普通 \instruction{mov} 在 \code{std::uint8_t} 返回值中有什么区别？
  \item \instruction{movss} 和 \instruction{movsd} 分别对应什么标量浮点类型？
  \item \code{load_2d} 的行长 5 如何进入汇编？
  \item 修改 \code{Pair} 字段顺序或添加字段后，\code{sizeof(Pair)}、字段 offset 和汇编如何变化？
\end{enumerate}
\end{labbox}

\subsection{实验 3.5：AoS 与 SoA 的第一次 benchmark}

\begin{labbox}
实现两个版本，计算所有粒子 \code{x} 坐标之和。

\begin{lstlisting}[language=C++]
struct Particle {
    float x;
    float y;
    float z;
    float mass;
};

float sum_x_aos(Particle const* p, std::size_t n)
{
    float s = 0.0f;
    for (std::size_t i = 0; i < n; ++i) {
        s += p[i].x;
    }
    return s;
}

float sum_x_soa(float const* x, std::size_t n)
{
    float s = 0.0f;
    for (std::size_t i = 0; i < n; ++i) {
        s += x[i];
    }
    return s;
}
\end{lstlisting}

要求：

\begin{enumerate}
  \item 使用相同的 \code{x} 数据初始化 AoS 和 SoA。
  \item 至少测试 \code{n = 1024}、\code{n = 1048576}、\code{n = 67108864}。
  \item 每个规模重复运行至少 20 次。
  \item 防止结果被编译器优化掉。
  \item 输出最小值、中位数、最大值。
  \item 查看 \code{-O3} 汇编，记录是否自动向量化。
\end{enumerate}

报告必须回答：

\begin{enumerate}
  \item 哪个版本更快？是否所有规模都一样？
  \item 从 cache line 利用率解释结果。
  \item 从汇编解释编译器是否生成了不同访问形态。
  \item 如果同时使用 \code{x,y,z,mass}，你的结论是否可能改变？
\end{enumerate}
\end{labbox}

\subsection{本章作业}

\begin{exercisebox}
\subsubsection*{作业 1：字节级对象报告}

选择 5 个对象：

\begin{itemize}
  \item 一个 \code{std::uint32_t}。
  \item 一个负的 \code{std::int32_t}。
  \item 一个 \code{float}。
  \item 一个含 padding 的结构体。
  \item 一个你自己设计的嵌套结构体。
\end{itemize}

对每个对象提交：

\begin{enumerate}
  \item 源码定义。
  \item \code{sizeof} 和 \code{alignof}。
  \item 字节 dump。
  \item 手工解释每个字段或位模式。
  \item 哪些字节不应该被当作稳定语义值。
\end{enumerate}

\subsubsection*{作业 2：结构体布局优化设计}

设计一个结构体表示简化版神经网络张量元数据，至少包含：

\begin{itemize}
  \item 数据指针。
  \item 维度数量。
  \item 4 个维度。
  \item 4 个 stride。
  \item dtype。
  \item flags。
  \item 引用计数或状态字段。
\end{itemize}

提交两个版本：

\begin{enumerate}
  \item 直觉字段顺序版本。
  \item 经过布局推导优化后的版本。
\end{enumerate}

要求给出每个版本的 \code{sizeof}、字段 offset、padding 分析，并解释你是为了减小大小、改善 cache 局部性，还是为了 ABI 可读性做取舍。

\subsubsection*{作业 3：指针和地址推导题}

假设：

\begin{lstlisting}[language=C++]
struct P {
    float x;
    float y;
    float z;
    float w;
};

P* p = reinterpret_cast<P*>(0x1000);
\end{lstlisting}

在 \code{sizeof(P) == 16} 的前提下，手工计算：

\begin{enumerate}
  \item \code{&p[0].x}
  \item \code{&p[0].z}
  \item \code{&p[3]}
  \item \code{&p[3].w}
  \item \code{reinterpret_cast<char*>(p) + 3}
\end{enumerate}

然后写程序在真实数组上验证相对偏移。报告中必须明确区分元素偏移和字节偏移。

\subsubsection*{作业 4：行为分类}

判断下列行为属于 well-defined、implementation-defined、unspecified 还是 undefined，并解释原因：

\begin{enumerate}
  \item 读取 \code{std::uint32_t} 的对象字节。
  \item 有符号 \code{int} 溢出。
  \item 无符号 \code{uint32_t} 溢出。
  \item 访问 \code{a[n]}，其中数组合法下标为 \code{0..n-1}。
  \item 用 \code{std::bit_cast<std::uint32_t>(float_value)} 读取位模式。
  \item 用 \code{reinterpret_cast<double*>(&int_obj)} 后解引用。
  \item 比较两个含 padding 的结构体对象的 \code{memcmp} 结果。
  \item 打印指针值并观察 ASLR 导致每次运行不同。
\end{enumerate}

\subsubsection*{作业 5：小型性能报告}

完成 AoS vs SoA benchmark。报告不少于 2000 字，必须包括：

\begin{itemize}
  \item 机器信息。
  \item 编译器和参数。
  \item 数据规模和初始化方式。
  \item correctness 验证方式。
  \item benchmark 方法。
  \item 汇编证据。
  \item 时间结果表。
  \item 用 cache line 利用率解释的性能模型。
  \item 至少一个失败或反直觉观察。
\end{itemize}
\end{exercisebox}

\subsection{验收标准}

完成本章后，你应该能做到：

\begin{itemize}
  \item 解释位、字节、对象表示、类型和值之间的关系。
  \item 读懂小端内存 dump。
  \item 区分无符号回绕、有符号补码表示和 C++ 有符号溢出 UB。
  \item 从 \code{a[i]} 推导出地址计算。
  \item 从 x86-64 内存操作数识别 base、index、scale、displacement。
  \item 手工推导简单结构体的 offset、padding、\code{sizeof} 和 \code{alignof}。
  \item 知道如何安全观察对象字节，什么时候应使用 \code{std::bit_cast} 或 \code{std::memcpy}。
  \item 解释别名和对象生命周期为什么影响编译器优化。
  \item 用 cache line 利用率初步解释连续访问、步长访问、随机访问、AoS 和 SoA 的性能差异。
\end{itemize}

如果你能做到这些，第 31 章的 cache、第 35 章的 AoS/SoA、第 37 章的 SIMD load/store、第 45 章的 GEMM blocking、第 50 章以后的 AI 算子布局优化，才会真正建立在可推导的基础上，而不是建立在经验口号上。
